\scanpage{88}

In the \unsure{last|second} place, let $f(s)$ be a function defined for $s>0$, such that $f(e^t)$ is convex, put $f(s_1,\mathord{-},s_n)=\sum_1^{\unsure{n}} f(s_i)$, it is a symmetric function of the \deleted{\unsure{numbers}}\inserted{\unsure{arguments}} of the $s_i>0$, and such that $f(e^{t_1},\mathord{-}e^{t_n})$ is convex. Consequently, \deleted{\unsure{$F$}} $f(s_1,\unsure{\mathord{-}|\deleted{+}}s_n)$ satisfies the conditions of corollary 1, under hypothesis b. If we further suppose $f(s)$ increasing, and defined and continuous also for $s\geq0$, then $f(s_1,\mathord{-}s_n)$ satisfies $=$ the conditions of \deleted{corollary 1}\inserted{\unsure{theorem|Theorem} 1, 2°} under hypothesis a. Hence, \deleted{\unsure{with the hypothesis}} \inserted{under the conditions $\unsure{\rho_i>0}$,} \deleted{\unsure{suitable on} $(\alpha_i),(\beta_i)$, \unsure{the inequalities}} the inequality

\[
\begin{array}{@{}l@{\qquad}l@{}}
(13)&\displaystyle \sum_1^n f(\rho_i)\leq\sum_1^n f(\sigma_i)
\end{array}
\]

and under the conditions of b), taking into account a preceding remark, the sequence of inequalities

\[
\begin{array}{@{}l@{\qquad}l@{}}
(14)&\displaystyle \sum_1^k f(\rho_i)\leq\sum_1^k f(\sigma_i)\qquad(k=1,\mathord{-}n)
\end{array}
\]

Let us take for example, \deleted{\unsure{$p$ positive and}}\inserted{for any real $p$} $f(s)=s^p$, \inserted{$(s>0)$} we then have $f(e^t)=e^{pt}$, which is indeed a convex function. \deleted{\unsure{$f$ itself}}\inserted{$f(r)$} is defined and \deleted{\unsure{even}} continuous for $s\geq0$, and \deleted{\unsure{an}} increasing function of $s$, as soon as $p\geq0$. \deleted{\unsure{the two} \illegible{1}}

{\reversemarginpar\margin{89}}
